% =====================================================================
%  APÊNDICE A — REGISTRO DE SOFTWARE NO INPI
% =====================================================================

\chapter{REGISTRO DE SOFTWARE NO INPI}
\label{ap:registro-inpi}

Este apêndice reúne os dados do registro do \argus{} junto ao Instituto
Nacional da Propriedade Industrial (INPI) e explicita o papel da versão
registrada como linha de base do projeto. O certificado oficial é reproduzido
nas páginas pré-textuais; a evolução técnica entre a versão registrada e a
versão apresentada neste trabalho é discutida no
Capítulo~\ref{cap:metodologia}.

\section{Contexto legal}
\label{ap:registro-contexto}

O registro de programa de computador é um direito de propriedade intelectual
previsto na Lei nº 9.609/1998 \cite{lei9609}. No Brasil, o registro é realizado
no INPI e assegura ao titular e aos autores a proteção jurídica sobre o código
fonte do software por 50 anos, contados de 1º de janeiro do ano subsequente à
sua publicação \cite{inpi2026}. O registro não é condição para a proteção
autoral, que decorre da criação da obra, mas constitui prova documental útil
para fins acadêmicos, comerciais e institucionais.

O \argus{} foi registrado com base na versão inicial do sistema, desenvolvida
na etapa inicial do projeto. Essa versão é tratada neste trabalho como
\textbf{linha de base} (\textit{baseline}) funcional e arquitetural. O registro
não congela o software: ele documenta um estado do código, enquanto o projeto
continuou a evoluir por meio de ciclos iterativos de desenvolvimento,
conforme descrito no Capítulo~\ref{cap:metodologia}.

\section{Dados do registro}
\label{ap:registro-dados}

A Tabela~\ref{tab:registro} consolida os dados do certificado.

\begin{table}[H]
    \centering
    \caption{Dados do registro do software no INPI}
    \label{tab:registro}
    \small
    \begin{tabular}{>{\raggedright\arraybackslash}p{4.6cm}p{9.2cm}}
        \toprule
        \textbf{Campo} & \textbf{Valor} \\
        \midrule
        Número do processo & BR512026003786-7 \\
        Título & Argus Cyclist \\
        Data de criação / publicação & 20/01/2026 \\
        Titular & Fundação Universidade do Estado do Rio Grande do Norte (FUERN) \\
        Autores & Isaac de Lima Oliveira Filho; Paulo Sérgio Silva de Medeiros \\
        Linguagem & HTML; JavaScript; Outros \\
        Campo de aplicação & IF-10 \\
        Tipo de programa & SM-02 \\
        Algoritmo \textit{hash} & SHA-512 \\
        Data de expedição & 02/06/2026 \\
        \bottomrule
    \end{tabular}
    \vspace{0.2cm}
    \par\small Fonte: certificado de registro do INPI (2026).
\end{table}

\section{A versão registrada como linha de base}
\label{ap:registro-versao}

A versão registrada contemplava as funcionalidades que comprovaram a
viabilidade técnica da proposta e definiram a identidade do projeto:

\begin{itemize}
    \item \textbf{Simulador núcleo}: conexão BLE com sensores de frequência
        cardíaca e rolos de treino inteligentes (FTMS), renderização de mapa
        3D e exportação de arquivos \texttt{.FIT}.
    \item \textbf{Histórico e gráficos}: banco de dados local (SQLite) com
        histórico de atividades e visualização gráfica de sessões.
    \item \textbf{Curva de potência (MMP)}: cálculo e exibição dos recordes
        de potência máxima em diferentes janelas de tempo.
    \item \textbf{Editor de rotas}: criação de rotas personalizadas
        diretamente no mapa.
\end{itemize}

Essa versão constituiu o \textit{baseline} funcional e arquitetural que foi
aprimorado ao longo do ciclo de vida do projeto. A Tabela~\ref{tab:evolucao}
do Capítulo~\ref{cap:metodologia} compara a versão registrada com a versão
atual, e os Capítulos~\ref{cap:arquitetura} e~\ref{cap:implementacao}
detalham as decisões e a implementação que produziram essa evolução.
